\section{Defect mass of the normalized band-product vector}
\label{sec:regional_low_defect_product_mass}

\begin{theorem}[Mass and nonnegative energy deficit]
\label{thm:regional_low_defect_product_mass}
\lean{Matrix.replicaMetric_normalized_symProj_bandProduct_lowDefect_mass_ge}
Let the physical and two auxiliary coordinate spaces be finite and nonzero.
Suppose that $\|\Omega\|=1$, $\gamma>0$, and
\begin{align*}
 H-E_0I\geq\gamma(I-|\Omega\rangle\langle\Omega|).
\end{align*}
Fix $k>0$, an isometry $Z$ onto simultaneous copy symmetry, and a nonzero
symmetric vector $\xi$ in the original auxiliary sectors
$\ell_C,\ell_R$. For terminal labels $j$ and bands $0\leq g<G$, let
$Q_{jg},Y_{jg}$ be disjoint physical regions such that, whenever $g<h$,
\begin{align*}
 Q_{jg}\cup Y_{jg}\subseteq Q_{j'h}
\end{align*}
for every pair of terminal labels $j,j'$. Let $t\geq0$, retain the
original full-system replica weights, and set
\begin{align*}
 A_{jg}&=\bigl(W(CQ_{jg})^{-1}
 W(R(Q_{jg}\cup Y_{jg})^c)^{-1}W(Y_{jg})\bigr)^2,\\
 B_{jg}&=Z^*A_{jg}Z,
 &M&=\prod_{g=0}^{G-1}T(B_{\cdot g}),
\end{align*}
where every band uses the same weighted mean tree $T$.
The symbols $C$ and $R$ in these regions denote the two auxiliary sites;
the complement in the second factor is taken among physical sites.
Let $\overline H_k$ be the physical replica average with both auxiliary
identities, in the global coordinates of
Definition~\ref{def:global_replica_mean_hamiltonian}.
For any real $s$, define
\begin{align*}
 w=M^{-s}Z^*\xi,\qquad
 \phi=\frac{w}{\|w\|},\qquad v=Z\phi,\qquad
 \delta=\frac{\operatorname{Re}\langle v,\overline H_kv\rangle-E_0}{\gamma}.
\end{align*}
Then $M$ is positive definite, $w\ne0$, $\|\phi\|=\|v\|=1$, and
$\delta\geq0$. For every $\tau>0$, the common projection
$\Pi_\tau$ from Definition~\ref{def:global_replica_low_defect_projection}
has an orthogonal projection $P_\tau=Z^*\Pi_\tau Z$ as its compression, and
\begin{align*}
 \operatorname{Re}\langle\phi,P_\tau\phi\rangle
 \geq 1-\frac{\delta}{\tau}.
\end{align*}
The cutoff parameter is chosen after the normalized vector. Repeated
terminal labels, zero tree weights, and $G=0$ are allowed.
This is the mass estimate used in
\cite[Section~8, lines~421--443]{OpenAI2026AreaLaw}.
\end{theorem}
\begin{proof}
\uses{thm:replica_metric_lifted_region_geometry,thm:replica_symmetric_band_label_preservation,thm:global_low_defect_gap_mass,thm:regional_low_defect_projection_global,def:global_replica_mean_hamiltonian}
The actual lifted regions satisfy the hypotheses of the label-preservation
theorem. Hence the same normalized vector $v$ has simultaneous symmetry
and both original auxiliary labels. Apply the original gap estimate to
$v$ for each $\tau>0$. The equality $v=Z\phi$ identifies its projection
mass with the displayed compressed quadratic form. Since the common
projection is contained in the symmetric range, its compression is an
orthogonal projection. Its mass on the unit vector $\phi$ is at most one.
Taking $\tau=1$ in the gap estimate therefore gives $1-\delta\leq1$,
and hence $\delta\geq0$.
\end{proof}
